\documentclass[10pt,aspectratio=169]{beamer}

% ====================== Tema y apariencia ======================
\usetheme{Madrid}
\usecolortheme{whale}
\useinnertheme{rounded}
\useoutertheme{infolines}
\setbeamertemplate{navigation symbols}{}
\setbeamertemplate{caption}[numbered]

% Pie de página mínimo: solo el número de diapositiva
\setbeamertemplate{footline}{%
  \hfill
  \usebeamercolor[fg]{page number in head/foot}%
  \usebeamerfont{page number in head/foot}%
  \footnotesize\insertframenumber\,/\,\inserttotalframenumber\kern1em\vskip3pt%
}

% ====================== Paquetes ======================
\usepackage[utf8]{inputenc}
\usepackage[T1]{fontenc}
\usepackage{lmodern}     % fuentes vectoriales (evita bitmaps; habilita microtype)
\usepackage{microtype}  % pulido tipográfico
\usepackage[spanish,es-noshorthands]{babel}
\usepackage{amsmath,amssymb}
\usepackage{booktabs}
\usepackage{listings}
\usepackage{pgfplots}
\pgfplotsset{compat=1.18}
\usepackage{tikz}
\usetikzlibrary{arrows.meta,positioning}
\hypersetup{unicode,  % marcadores del PDF (índice) con acentos y ¿ ¡ correctos
  pdftitle={Cómo crear documentos y presentaciones con LaTeX},
  pdfauthor={Rodrigo Martín Vidal Galán}}

% Colores propios
\definecolor{azulPrincipal}{RGB}{20,60,120}
\definecolor{verdeAgua}{RGB}{0,140,150}
\definecolor{naranja}{RGB}{210,110,30}
\definecolor{vino}{RGB}{170,30,70}

\setbeamercolor{structure}{fg=azulPrincipal}
\setbeamercolor{frametitle}{bg=azulPrincipal,fg=white}

% Comando propio de ejemplo (se muestra en el nivel avanzado)
\newcommand{\Rdemo}{\mathbb{R}}

% Estilo para mostrar código LaTeX
\lstset{
  language=[LaTeX]TeX,
  basicstyle=\ttfamily\scriptsize,
  keywordstyle=\color{azulPrincipal}\bfseries,
  commentstyle=\color{gray}\itshape,
  stringstyle=\color{naranja},
  backgroundcolor=\color{azulPrincipal!6},
  frame=single,
  rulecolor=\color{azulPrincipal!30},
  breaklines=true,
  showstringspaces=false,
  columns=fullflexible,
  keepspaces=true,
  texcsstyle=*\color{azulPrincipal}\bfseries,
  extendedchars=true,
  literate=
    {á}{{\'a}}1 {é}{{\'e}}1 {í}{{\'i}}1 {ó}{{\'o}}1 {ú}{{\'u}}1
    {Á}{{\'A}}1 {É}{{\'E}}1 {Í}{{\'I}}1 {Ó}{{\'O}}1 {Ú}{{\'U}}1
    {ñ}{{\~n}}1 {Ñ}{{\~N}}1 {¿}{{?`}}1 {¡}{{!`}}1 {°}{{\textdegree}}1,
}

% Divisoria de nivel
\newcommand{\nivel}[3]{%
  \begin{frame}[plain]
    \centering\vfill
    {\small\structure{#1}}\\[8pt]
    {\Large\textbf{#2}}\\[12pt]
    {\normalsize #3}
    \vfill
  \end{frame}}

% ====================== Datos del título ======================
\title[LaTeX y Beamer]{Cómo crear documentos y presentaciones con LaTeX}
\subtitle{Una introducción práctica desde cero (¡hecha en LaTeX!)}
\author{Rodrigo Martín Vidal Galán}
\institute{}
\date{}

\begin{document}

% ====================== Portada ======================
{
\setbeamertemplate{footline}{}
\begin{frame}[plain]
  \titlepage
\end{frame}
}

% ====================== Agenda ======================
\begin{frame}{Recorrido por niveles}
  \tableofcontents
  \vfill
  \begin{center}\small De lo \textbf{básico} a lo \textbf{avanzado}, y cerramos con \textbf{ejemplos integrados}.\end{center}
\end{frame}

% ####################################################################
\section{Nivel 1 -- Lo básico}
% ####################################################################
\nivel{NIVEL 1}{Lo básico}{Tu primer documento: qué es LaTeX, cómo se compila, estructura, texto, listas y matemática elemental.}

\begin{frame}{¿Qué es LaTeX?}
  \begin{block}{Idea central}
    LaTeX es un \textbf{sistema de composición de documentos}: vos escribís
    texto plano con \emph{instrucciones}, y LaTeX se encarga de la
    \textbf{apariencia} (tipografía, espaciado, fórmulas, figuras).
  \end{block}
  \begin{columns}[T]
    \begin{column}{0.48\textwidth}
      \textbf{Procesador tradicional (Word)}
      \begin{itemize}
        \item Ves el resultado mientras escribís (WYSIWYG).
        \item Vos ajustás el formato a mano.
      \end{itemize}
    \end{column}
    \hspace{0.04\textwidth}
    \begin{column}{0.48\textwidth}
      \textbf{LaTeX}
      \begin{itemize}
        \item Escribís el \emph{contenido}; el formato es automático.
        \item Ideal para \textbf{matemática}, tesis, libros y presentaciones.
        \item Resultados consistentes y profesionales.
      \end{itemize}
    \end{column}
  \end{columns}
\end{frame}

\begin{frame}{Herramientas y flujo de trabajo}
  \begin{columns}[T]
    \begin{column}{0.47\textwidth}
      \begin{block}{¿Qué necesito?}
        \begin{itemize}
          \item Una \textbf{distribución} LaTeX: \texttt{MiKTeX} (Windows) o
            \texttt{TeX~Live}.
          \item Un \textbf{editor}: TeXstudio, VS Code, o \texttt{Overleaf}
            (online, sin instalar nada).
        \end{itemize}
      \end{block}
    \end{column}
    \hspace{0.04\textwidth}
    \begin{column}{0.47\textwidth}
      \begin{block}{El ciclo de trabajo}
        \[ \texttt{archivo.tex} \;\xrightarrow{\text{compilar}}\; \texttt{archivo.pdf} \]
        \begin{enumerate}
          \item Escribís el \texttt{.tex}.
          \item \textbf{Compilás} (con \texttt{pdflatex} o \texttt{latexmk}).
          \item Mirás el \texttt{.pdf} y repetís.
        \end{enumerate}
      \end{block}
    \end{column}
  \end{columns}
\end{frame}

\begin{frame}[fragile]{Instalación paso a paso (Windows)}
  \begin{columns}[T]
    \begin{column}{0.52\textwidth}
      \begin{enumerate}
        \item Descargá e instalá \textbf{MiKTeX} desde \texttt{miktex.org}.
        \item Instalá un editor: \textbf{TeXstudio} o \textbf{VS~Code} con la
          extensión \emph{LaTeX Workshop}.
        \item Creá un archivo \texttt{.tex} y escribí tu documento.
        \item Apretá \textbf{Compilar} (o usá la consola).
      \end{enumerate}
    \end{column}
    \hspace{0.03\textwidth}
    \begin{column}{0.44\textwidth}
      \begin{block}{Compilar desde la consola}
\begin{lstlisting}
pdflatex documento.tex
% o, recomendado (resuelve
% referencias en varias pasadas):
latexmk -pdf documento.tex
\end{lstlisting}
      \end{block}
      \begin{exampleblock}{Tip}
        MiKTeX instala los paquetes que falten \textbf{automáticamente} la
        primera vez que los usás.
      \end{exampleblock}
    \end{column}
  \end{columns}
\end{frame}

\begin{frame}[fragile]{La estructura mínima de un .tex}
  \begin{columns}[T]
    \begin{column}{0.55\textwidth}
\begin{lstlisting}
\documentclass{article}

% --- Preámbulo ---
\usepackage[utf8]{inputenc}
\title{Mi documento}
\author{Yo}

% --- Cuerpo ---
\begin{document}
\maketitle
Hola, mundo.
\end{document}
\end{lstlisting}
    \end{column}
    \hspace{0.03\textwidth}
    \begin{column}{0.41\textwidth}
      \begin{itemize}
        \item \texttt{\textbackslash documentclass}: el tipo de documento.
        \item \textbf{Preámbulo}: todo lo que va \emph{antes} de
          \texttt{\textbackslash begin\{document\}} (paquetes, ajustes).
        \item \textbf{Cuerpo}: entre \texttt{\textbackslash begin\{document\}}
          y \texttt{\textbackslash end\{document\}}.
        \item Las líneas con \texttt{\%} son \textbf{comentarios}.
      \end{itemize}
    \end{column}
  \end{columns}
\end{frame}

\begin{frame}[fragile]{Comandos, entornos y paquetes}
  \begin{columns}[c]
    \begin{column}{0.52\textwidth}
      \begin{block}{Dos ladrillos fundamentales}
        \textbf{Comando:} empieza con \texttt{\textbackslash} y puede llevar
        argumentos entre llaves.
\begin{lstlisting}
\textbf{texto en negrita}
\end{lstlisting}
        \textbf{Entorno:} bloque con apertura y cierre.
\begin{lstlisting}
\begin{center}
  texto centrado
\end{center}
\end{lstlisting}
      \end{block}
    \end{column}
    \hspace{0.03\textwidth}
    \begin{column}{0.45\textwidth}
      \begin{block}{Paquetes}
        Se cargan en el preámbulo y agregan funciones:
\begin{lstlisting}
\usepackage{amsmath} % matematica
\usepackage{graphicx}% imagenes
\usepackage{tikz}    % graficos
\end{lstlisting}
      \end{block}
    \end{column}
  \end{columns}
\end{frame}

\begin{frame}[fragile]{Formato de texto}
  \begin{columns}[c]
    \begin{column}{0.54\textwidth}
\begin{lstlisting}
Texto normal,
\textbf{negrita},
\emph{énfasis (itálica)} y
\texttt{monoespaciado}.
\end{lstlisting}
    \end{column}
    \hspace{0.04\textwidth}
    \begin{column}{0.42\textwidth}
      \begin{exampleblock}{Resultado}
        Texto normal, \textbf{negrita}, \emph{énfasis (itálica)} y
        \texttt{monoespaciado}.
      \end{exampleblock}
    \end{column}
  \end{columns}
\end{frame}

\begin{frame}[fragile]{Listas}
  \begin{columns}[c]
    \begin{column}{0.54\textwidth}
\begin{lstlisting}
\begin{itemize}
  \item Primer punto
  \item Segundo punto
\end{itemize}

\begin{enumerate}
  \item Paso uno
  \item Paso dos
\end{enumerate}
\end{lstlisting}
    \end{column}
    \hspace{0.04\textwidth}
    \begin{column}{0.42\textwidth}
      \begin{exampleblock}{Resultado}
        \begin{itemize}
          \item Primer punto
          \item Segundo punto
        \end{itemize}
        \begin{enumerate}
          \item Paso uno
          \item Paso dos
        \end{enumerate}
      \end{exampleblock}
    \end{column}
  \end{columns}
\end{frame}

\begin{frame}[fragile]{Matemática: en línea y centrada}
  \begin{columns}[c]
    \begin{column}{0.54\textwidth}
\begin{lstlisting}
En línea: $a^2 + b^2 = c^2$.

Centrada (display):
\[ x = \frac{-b \pm \sqrt{b^2-4ac}}{2a} \]
\end{lstlisting}
      \scriptsize El modo matemático va entre \texttt{\$...\$} (en línea)
      o \texttt{\textbackslash[...\textbackslash]} (centrado).
    \end{column}
    \hspace{0.04\textwidth}
    \begin{column}{0.42\textwidth}
      \begin{exampleblock}{Resultado}
        En línea: $a^2+b^2=c^2$.
        \medskip

        Centrada (display):
        \[ x=\frac{-b\pm\sqrt{b^2-4ac}}{2a} \]
      \end{exampleblock}
    \end{column}
  \end{columns}
\end{frame}

\begin{frame}[fragile]{Matemática: símbolos más usados}
  \begin{columns}[c]
    \begin{column}{0.54\textwidth}
\begin{lstlisting}
$x_1$, $x^2$        % sub/super
$\frac{1}{2}$       % fracción
$\sqrt{x}$          % raíz
$\sum_{i=1}^{n} i$  % sumatoria
$\alpha,\ \pi,\ \infty$
$\leq,\ \geq,\ \neq$
\end{lstlisting}
      \scriptsize Muchos símbolos vienen del paquete \texttt{amsmath}.
    \end{column}
    \hspace{0.04\textwidth}
    \begin{column}{0.42\textwidth}
      \begin{exampleblock}{Resultado}
        $x_1$,\quad $x^2$ \\[3pt]
        $\dfrac{1}{2}$,\quad $\sqrt{x}$,\quad
        $\displaystyle\sum_{i=1}^{n} i$ \\[6pt]
        $\alpha,\ \pi,\ \infty$ \\[3pt]
        $\leq,\ \geq,\ \neq$
      \end{exampleblock}
    \end{column}
  \end{columns}
\end{frame}

% ####################################################################
\section{Nivel 2 -- Intermedio}
% ####################################################################
\nivel{NIVEL 2}{Intermedio}{Documentos más ricos y tu primera presentación: tablas, imágenes, referencias, más matemática, Beamer y un primer gráfico.}

\begin{frame}[fragile]{Tablas}
  \begin{columns}[c]
    \begin{column}{0.54\textwidth}
\begin{lstlisting}
\begin{tabular}{l c}
  \toprule
  Forma & Vértice \\
  \midrule
  Canónica & (h, k) \\
  Polinómica & (0, c) \\
  \bottomrule
\end{tabular}
\end{lstlisting}
      \scriptsize (\texttt{\textbackslash toprule}, etc. son del paquete
      \texttt{booktabs}.)
    \end{column}
    \hspace{0.04\textwidth}
    \begin{column}{0.42\textwidth}
      \begin{exampleblock}{Resultado}
        \centering
        \begin{tabular}{l c}
          \toprule
          Forma & Vértice \\
          \midrule
          Canónica & $(h,k)$ \\
          Polinómica & $(0,c)$ \\
          \bottomrule
        \end{tabular}
      \end{exampleblock}
    \end{column}
  \end{columns}
\end{frame}

\begin{frame}[fragile]{Imágenes y figuras}
  \begin{columns}[c]
    \begin{column}{0.54\textwidth}
\begin{lstlisting}
\usepackage{graphicx} % preámbulo

\begin{figure}
  \centering
  \includegraphics[width=0.6\linewidth]
                  {foto.png}
  \caption{Una parábola.}
\end{figure}
\end{lstlisting}
      \scriptsize \texttt{figure} la posiciona y le agrega numeración con
      \texttt{\textbackslash caption}.
    \end{column}
    \hspace{0.04\textwidth}
    \begin{column}{0.42\textwidth}
      \begin{exampleblock}{Resultado}
        \centering
        \begin{tikzpicture}
          \draw[gray!60] (0,0) rectangle (3.4,2.0);
          \begin{scope}[shift={(1.7,0.3)}]
            \draw[azulPrincipal,thick,smooth,domain=-1.15:1.15,samples=40]
              plot ({\x},{0.95*\x*\x});
          \end{scope}
        \end{tikzpicture}\\[2pt]
        {\footnotesize Figura 1: Una parábola.}
      \end{exampleblock}
    \end{column}
  \end{columns}
\end{frame}

\begin{frame}[fragile]{Referencias cruzadas y notas al pie}
  \begin{columns}[c]
    \begin{column}{0.54\textwidth}
\begin{lstlisting}
La ecuación \eqref{eq:res} da las raíces.

\begin{equation}\label{eq:res}
  x=\frac{-b\pm\sqrt{b^2-4ac}}{2a}
\end{equation}

Una aclaración.\footnote{Texto al pie.}
\end{lstlisting}
      \scriptsize \texttt{\textbackslash label} marca y
      \texttt{\textbackslash ref}/\texttt{\textbackslash eqref} recuperan el
      número: si reordenás, se actualiza solo.
    \end{column}
    \hspace{0.04\textwidth}
    \begin{column}{0.42\textwidth}
      \begin{exampleblock}{Resultado}
        La ecuación~(1) da las raíces.
        \[ x=\frac{-b\pm\sqrt{b^2-4ac}}{2a} \tag{1} \]
        Una aclaración.$^{1}$
        \vspace{4pt}\hrule\vspace{2pt}
        {\scriptsize $^{1}$ Texto al pie.}
      \end{exampleblock}
    \end{column}
  \end{columns}
\end{frame}

\begin{frame}[fragile]{Matemática: matrices y sistemas}
  \begin{columns}[c]
    \begin{column}{0.54\textwidth}
\begin{lstlisting}
A = \begin{pmatrix}
  1 & 2 \\
  3 & 4
\end{pmatrix}

f(x)=\begin{cases}
  x^2 & \text{si } x\ge 0\\
  -x  & \text{si } x<0
\end{cases}
\end{lstlisting}
      \scriptsize \texttt{pmatrix} y \texttt{cases} vienen de \texttt{amsmath}.
    \end{column}
    \hspace{0.04\textwidth}
    \begin{column}{0.42\textwidth}
      \begin{exampleblock}{Resultado}
        \[ A=\begin{pmatrix} 1 & 2 \\ 3 & 4 \end{pmatrix} \]
        \[ f(x)=\begin{cases} x^2 & \text{si } x\ge 0\\ -x & \text{si } x<0 \end{cases} \]
      \end{exampleblock}
    \end{column}
  \end{columns}
\end{frame}

\begin{frame}[fragile]{Matemática: alinear ecuaciones, integrales y límites}
  \begin{columns}[c]
    \begin{column}{0.54\textwidth}
\begin{lstlisting}
\begin{align*}
  (x+1)^2 &= x^2 + 2x + 1 \\
          &= x^2 + 2x + 1
\end{align*}

$\int_0^1 x^2\,dx = \tfrac{1}{3}$

$\lim_{x\to 0}\frac{\sin x}{x}=1$
\end{lstlisting}
      \scriptsize En \texttt{align} el \texttt{\&} marca por dónde se alinean
      las ecuaciones (acá, el \texttt{=}).
    \end{column}
    \hspace{0.04\textwidth}
    \begin{column}{0.42\textwidth}
      \begin{exampleblock}{Resultado}
        \begin{align*}
          (x+1)^2 &= x^2 + 2x + 1 \\
                  &= x^2 + 2x + 1
        \end{align*}
        \[ \int_0^1 x^2\,dx=\tfrac{1}{3} \]
        \[ \lim_{x\to 0}\frac{\sin x}{x}=1 \]
      \end{exampleblock}
    \end{column}
  \end{columns}
\end{frame}

\begin{frame}{\texttt{article} vs Beamer: ¿cuál elijo?}
  Las dos son \textbf{clases de documento}: cambia el
  \texttt{\textbackslash documentclass}.
  \medskip
  \renewcommand{\arraystretch}{1.3}
  \begin{center}
  \footnotesize
  \begin{tabular}{@{}p{0.18\textwidth} p{0.35\textwidth} p{0.35\textwidth}@{}}
    \toprule
    \textbf{Aspecto} & \texttt{\textbf{article}} & \texttt{\textbf{beamer}} \\
    \midrule
    ¿Para qué? & Documentos para leer o imprimir: apuntes, informes, TPs, tesis. & Presentaciones para proyectar: diapositivas. \\
    Unidad básica & La \textbf{página} (el texto fluye). & La \textbf{diapositiva} (\texttt{frame}). \\
    Se arma con & \texttt{\textbackslash section}, párrafos, figuras. & \texttt{frame}, bloques, columnas. \\
    Portada & \texttt{\textbackslash maketitle}. & \texttt{\textbackslash titlepage}. \\
    \bottomrule
  \end{tabular}
  \end{center}
  \vfill
  \begin{exampleblock}{Lo que comparten}
    El \textbf{mismo lenguaje}: preámbulo, paquetes, matemática, tablas,
    imágenes y gráficos funcionan igual en ambas. ¡Lo que aprendés sirve para
    las dos!
  \end{exampleblock}
\end{frame}

\begin{frame}[fragile]{Beamer: cada \texttt{frame} es una diapositiva}
  \begin{columns}[c]
    \begin{column}{0.54\textwidth}
      Se usa \texttt{\textbackslash documentclass\{beamer\}}. Cada diapositiva
      es un entorno \texttt{frame}:
      \lstinputlisting{ejemplo_frame_beamer.tex}
    \end{column}
    \hspace{0.04\textwidth}
    \begin{column}{0.42\textwidth}
      \begin{exampleblock}{Resultado (una diapositiva)}
        \centering
        \fbox{\begin{minipage}{0.9\linewidth}
          \color{azulPrincipal}\rule{\linewidth}{6pt}\\[2pt]
          \footnotesize\textbf{Mi primera diapositiva}\\[4pt]
          \color{black}\footnotesize Hola, soy el contenido de la diapositiva.
          \vspace{4pt}
        \end{minipage}}
      \end{exampleblock}
    \end{column}
  \end{columns}
\end{frame}

\begin{frame}[fragile]{Beamer: bloques de color}
  \begin{columns}[c]
    \begin{column}{0.5\textwidth}
\begin{lstlisting}
\begin{block}{Definición}
  Contenido del bloque.
\end{block}

\begin{alertblock}{¡Atención!}
  Algo importante.
\end{alertblock}

\begin{exampleblock}{Ejemplo}
  Un tercer color.
\end{exampleblock}
\end{lstlisting}
    \end{column}
    \hspace{0.04\textwidth}
    \begin{column}{0.46\textwidth}
      \begin{block}{Definición}
        Contenido del bloque.
      \end{block}
      \begin{alertblock}{¡Atención!}
        Algo importante.
      \end{alertblock}
      \begin{exampleblock}{Ejemplo}
        Un tercer color.
      \end{exampleblock}
    \end{column}
  \end{columns}
\end{frame}

\begin{frame}[fragile]{Beamer: dos columnas}
  Sirven para poner contenido \textbf{lado a lado} (texto y gráfico, o
  código y resultado, como en toda esta presentación):
  \begin{columns}[c]
    \begin{column}{0.54\textwidth}
\begin{lstlisting}
\begin{columns}
  \begin{column}{0.5\textwidth}
    Izquierda
  \end{column}
  \begin{column}{0.5\textwidth}
    Derecha
  \end{column}
\end{columns}
\end{lstlisting}
    \end{column}
    \hspace{0.04\textwidth}
    \begin{column}{0.42\textwidth}
      \begin{exampleblock}{Resultado}
        \fbox{\begin{minipage}[t]{0.43\linewidth}\centering\strut Izquierda\end{minipage}}%
        \hfill
        \fbox{\begin{minipage}[t]{0.43\linewidth}\centering\strut Derecha\end{minipage}}
      \end{exampleblock}
    \end{column}
  \end{columns}
\end{frame}

\begin{frame}[fragile]{Beamer: temas y colores}
  \begin{columns}[c]
    \begin{column}{0.5\textwidth}
      Una o dos líneas en el \textbf{preámbulo} cambian todo el diseño
      (disposición y colores):
\begin{lstlisting}
\usetheme{Madrid}
\usecolortheme{whale}
\end{lstlisting}
      Hay muchos temas: \texttt{Berlin}, \texttt{Singapore},
      \texttt{Copenhagen}, \texttt{metropolis}\dots
    \end{column}
    \hspace{0.04\textwidth}
    \begin{column}{0.46\textwidth}
      \begin{exampleblock}{El mismo título, distintos colores}
        \footnotesize
        \colorbox{azulPrincipal}{\makebox[0.9\linewidth][l]{\color{white}\strut\ Mi diapositiva}}\\[4pt]
        \colorbox{vino}{\makebox[0.9\linewidth][l]{\color{white}\strut\ Mi diapositiva}}\\[4pt]
        \colorbox{verdeAgua}{\makebox[0.9\linewidth][l]{\color{white}\strut\ Mi diapositiva}}
        \\[4pt]
        \scriptsize Cambia la barra de título, las fuentes y los colores.
      \end{exampleblock}
    \end{column}
  \end{columns}
\end{frame}

\begin{frame}[fragile]{Gráficos: una parábola con \texttt{pgfplots}}
  \begin{columns}[c]
    \begin{column}{0.54\textwidth}
\begin{lstlisting}
\begin{tikzpicture}
\begin{axis}[width=\linewidth, height=4.5cm,
    axis lines=middle, xlabel=$x$, ylabel=$y$,
    xmin=-3.2, xmax=3.2, ymin=-0.7, ymax=9.5,
    xtick={-2,2}, ytick={4,8}]
  \addplot[blue,very thick,
           domain=-3:3,samples=80]{x^2};
\end{axis}
\end{tikzpicture}
\end{lstlisting}
    \end{column}
    \hspace{0.04\textwidth}
    \begin{column}{0.42\textwidth}
      \begin{exampleblock}{Resultado}
        \begin{tikzpicture}
        \begin{axis}[width=\linewidth, height=4.5cm,
            axis lines=middle, xlabel=$x$, ylabel=$y$,
            xmin=-3.2, xmax=3.2, ymin=-0.7, ymax=9.5,
            xtick={-2,2}, ytick={4,8}]
          \addplot[blue,very thick,
                   domain=-3:3,samples=80]{x^2};
        \end{axis}
        \end{tikzpicture}
      \end{exampleblock}
    \end{column}
  \end{columns}
\end{frame}

% ####################################################################
\section{Nivel 3 -- Avanzado}
% ####################################################################
\nivel{NIVEL 3}{Avanzado}{Personalización y herramientas potentes: comandos propios, overlays, gráficos avanzados, diagramas, bibliografía y errores comunes.}

\begin{frame}[fragile]{Comandos propios y colores}
  \begin{columns}[c]
    \begin{column}{0.54\textwidth}
\begin{lstlisting}
% definir un atajo:
\newcommand{\R}{\mathbb{R}}
$f : \R \to \R$

% definir y usar un color:
\definecolor{miazul}{RGB}{20,60,120}
\textcolor{miazul}{texto en mi azul}
\end{lstlisting}
      \scriptsize \texttt{\textbackslash newcommand} crea atajos para lo que
      repetís mucho; mantiene el documento prolijo y consistente.
    \end{column}
    \hspace{0.04\textwidth}
    \begin{column}{0.42\textwidth}
      \begin{exampleblock}{Resultado}
        $f : \Rdemo \to \Rdemo$
        \medskip

        \textcolor{azulPrincipal}{texto en mi azul}
      \end{exampleblock}
    \end{column}
  \end{columns}
\end{frame}

\begin{frame}[fragile]{Beamer avanzado: aparición por pasos}
  \begin{columns}[c]
    \begin{column}{0.54\textwidth}
\begin{lstlisting}
\begin{itemize}
  \item<1-> Aparece primero
  \item<2-> Aparece al segundo clic
\end{itemize}

Texto \pause y más texto después.
\end{lstlisting}
      \scriptsize Los \emph{overlays} (\texttt{<1->}, \texttt{\textbackslash pause})
      revelan el contenido de a poco. En \textbf{estas} presentaciones los
      evitamos para que sean lineales.
    \end{column}
    \hspace{0.04\textwidth}
    \begin{column}{0.42\textwidth}
      \begin{exampleblock}{Resultado (por clics)}
        \footnotesize
        \textbf{Clic 1:}
        \begin{itemize}
          \item Aparece primero
        \end{itemize}
        \textbf{Clic 2:}
        \begin{itemize}
          \item Aparece primero
          \item Aparece al segundo clic
        \end{itemize}
      \end{exampleblock}
    \end{column}
  \end{columns}
\end{frame}

\begin{frame}[fragile]{pgfplots avanzado: leyenda, varias curvas y puntos}
  \begin{columns}[c]
    \begin{column}{0.55\textwidth}
\begin{lstlisting}
\begin{tikzpicture}
\begin{axis}[width=\linewidth, height=4.5cm,
    axis lines=middle, xlabel=$x$,
    legend pos=north west]
  \addplot[blue,very thick,domain=-2:2]{x^2};
  \addlegendentry{$x^2$}
  \addplot[orange,very thick,domain=-2:2]{2*x+1};
  \addlegendentry{$2x+1$}
  \addplot[red,only marks,mark=*]
    coordinates {(1,1)};
\end{axis}
\end{tikzpicture}
\end{lstlisting}
    \end{column}
    \hspace{0.03\textwidth}
    \begin{column}{0.42\textwidth}
      \begin{exampleblock}{Resultado}
        \begin{tikzpicture}
        \begin{axis}[width=\linewidth, height=4.5cm,
            axis lines=middle, xlabel=$x$,
            legend pos=north west]
          \addplot[blue,very thick,domain=-2:2]{x^2};
          \addlegendentry{$x^2$}
          \addplot[orange,very thick,domain=-2:2]{2*x+1};
          \addlegendentry{$2x+1$}
          \addplot[red,only marks,mark=*]
            coordinates {(1,1)};
        \end{axis}
        \end{tikzpicture}
      \end{exampleblock}
    \end{column}
  \end{columns}
\end{frame}

\begin{frame}[fragile]{TikZ: diagramas con nodos y flechas}
  \begin{columns}[c]
    \begin{column}{0.54\textwidth}
\begin{lstlisting}
\usetikzlibrary{arrows.meta,
                positioning}
\begin{tikzpicture}
  \node (a) {.tex};
  \node[right=of a] (b) {compilar};
  \node[right=of b] (c) {.pdf};
  \draw[-{Stealth}] (a) -- (b);
  \draw[-{Stealth}] (b) -- (c);
\end{tikzpicture}
\end{lstlisting}
      \scriptsize TikZ sirve para dibujar: nodos, flechas, diagramas de flujo,
      esquemas\dots
    \end{column}
    \hspace{0.04\textwidth}
    \begin{column}{0.42\textwidth}
      \begin{exampleblock}{Resultado}
        \centering
        \begin{tikzpicture}[
            node distance=0.5cm and 0.7cm,
            every node/.style={font=\scriptsize},
            caja/.style={draw,rounded corners,fill=azulPrincipal!10,inner sep=4pt},
          ]
          \node[caja] (a) {.tex};
          \node[caja,right=of a] (b) {compilar};
          \node[caja,right=of b] (c) {.pdf};
          \draw[-{Stealth},azulPrincipal,thick] (a) -- (b);
          \draw[-{Stealth},azulPrincipal,thick] (b) -- (c);
        \end{tikzpicture}
      \end{exampleblock}
    \end{column}
  \end{columns}
\end{frame}

\begin{frame}[fragile]{Bibliografía y citas}
  \begin{columns}[c]
    \begin{column}{0.54\textwidth}
\begin{lstlisting}
% archivo "refs.bib":
@book{carena2019,
  author = {Carena, Marilina},
  title  = {Manual de matemática
            preuniversitaria},
  year   = {2019}
}
% en el documento:
Según \cite{carena2019} ...
\printbibliography
\end{lstlisting}
      \scriptsize Con \texttt{biblatex} la cita y la lista de referencias se
      arman y numeran \textbf{solas}.
    \end{column}
    \hspace{0.04\textwidth}
    \begin{column}{0.42\textwidth}
      \begin{exampleblock}{Resultado}
        \footnotesize
        Según [1] $\dots$
        \vspace{6pt}\hrule\vspace{4pt}
        \textbf{Referencias}\\[2pt]
        [1] Carena, M. (2019). \emph{Manual de matemática preuniversitaria}.
        Ediciones UNL.
      \end{exampleblock}
    \end{column}
  \end{columns}
\end{frame}

\begin{frame}{Errores comunes y cómo resolverlos}
  \small
  \renewcommand{\arraystretch}{1.35}
  \begin{tabular}{@{}p{0.34\textwidth} p{0.6\textwidth}@{}}
    \toprule
    \textbf{Mensaje / síntoma} & \textbf{Causa y solución} \\
    \midrule
    \texttt{Undefined control sequence} &
      Comando mal escrito o falta un \texttt{\textbackslash usepackage}.
      Revisá la ortografía y cargá el paquete. \\
    \texttt{Missing \$ inserted} &
      Un símbolo matemático fuera de modo math. Encerralo en \texttt{\$...\$}. \\
    \texttt{Runaway argument} / falta \texttt{\}} &
      Una llave o un entorno sin cerrar. Revisá \texttt{\textbackslash begin}/\texttt{\textbackslash end} y las llaves. \\
    Error con \texttt{verbatim}/código &
      Un \texttt{frame} con código sin la opción \texttt{[fragile]}. Agregala. \\
    \bottomrule
  \end{tabular}
  \vfill
  \begin{exampleblock}{Tip de oro}
    Leé \textbf{la primera} línea de error y el \textbf{número de línea}: casi
    siempre el problema está ahí (o una línea antes).
  \end{exampleblock}
\end{frame}

% ####################################################################
\section{Nivel 4 -- Ejemplos integrados}
% ####################################################################
\nivel{NIVEL 4}{Ejemplos integrados}{Casos reales que combinan varias herramientas a la vez: desde la carátula hasta un problema resuelto de punta a punta.}

\begin{frame}[fragile]{La carátula: el código}
  En el \textbf{preámbulo} declarás los datos y en el \textbf{cuerpo} los
  mostrás (\texttt{\textbackslash maketitle} en \texttt{article},
  \texttt{\textbackslash titlepage} en Beamer):
  \lstinputlisting{ejemplo_caratula.tex}
\end{frame}

\begin{frame}{La carátula: el resultado}
  \centering
  {\footnotesize\itshape (representación: la portada ocupa una página completa)}
  \vspace{10pt}

  \fbox{\begin{minipage}[t]{0.7\linewidth}
    \vspace{18pt}\centering
    \colorbox{azulPrincipal}{\makebox[0.9\linewidth][c]{\color{white}\large\strut Función cuadrática}}\\[18pt]
    {\normalsize Grupo 5}\\[4pt]
    {\normalsize Junio de 2026}
    \vspace{20pt}
  \end{minipage}}
\end{frame}

\begin{frame}[fragile]{Anatomía de una diapositiva: el código}
  Así se arma la diapositiva del ``máximo beneficio'' (Actividad~6): un
  \texttt{frame} con dos columnas, una fórmula, un \texttt{alertblock} y un
  gráfico. En la \textbf{siguiente} diapositiva, este mismo código renderizado
  de verdad:
  \lstinputlisting{ejemplo_anatomia.tex}
\end{frame}

\begin{frame}{Máximo beneficio}
  \begin{columns}
    \begin{column}{0.5\textwidth}
      $P(x)=5000+1000x-5x^2$, con $a<0$.
      \[ x=-\frac{b}{2a}=100 \]
      \begin{alertblock}{Conclusión}
        Máximo $55000$ en $x=100$.
      \end{alertblock}
    \end{column}
    \begin{column}{0.5\textwidth}
      \begin{tikzpicture}
      \begin{axis}[width=\linewidth,height=4cm,
                   axis lines=middle,xlabel=$x$]
        \addplot[blue,very thick,domain=0:200,
                 samples=40]{5000+1000*x-5*x^2};
        \addplot[red,only marks,mark=*]
          coordinates {(100,55000)};
      \end{axis}
      \end{tikzpicture}
    \end{column}
  \end{columns}
\end{frame}

\begin{frame}[fragile]{Un documento \texttt{article}: el código}
\begin{lstlisting}
\documentclass{article}
\usepackage[utf8]{inputenc}
\usepackage{amsmath}
\title{Informe de laboratorio}
\author{Grupo 5}
\date{\today}

\begin{document}
\maketitle
\section{Introducción}
Estudiamos la función cuadrática.
\section{Resultados}
La raíz se obtiene con
\begin{equation}
  x=\frac{-b\pm\sqrt{b^2-4ac}}{2a}
\end{equation}
\end{document}
\end{lstlisting}
\end{frame}

\begin{frame}{Un documento \texttt{article}: el resultado}
  \centering
  {\footnotesize\itshape (representación de la página generada)}
  \vspace{8pt}

  \fbox{\begin{minipage}[t]{0.78\linewidth}
    \centering
    {\large\textbf{Informe de laboratorio}}\\[2pt]
    {\small Grupo 5 \quad\today}\\[12pt]
    \raggedright\small
    \textbf{1\quad Introducción}\\[2pt]
    Estudiamos la función cuadrática.\\[10pt]
    \textbf{2\quad Resultados}\\[2pt]
    La raíz se obtiene con
    \[ x=\frac{-b\pm\sqrt{b^2-4ac}}{2a} \tag{1} \]
    \vspace{4pt}
  \end{minipage}}
\end{frame}

\begin{frame}[fragile]{Un problema resuelto: el código}
  \footnotesize Dentro de un \texttt{frame} combinamos texto, \texttt{align},
  un \texttt{block} y un gráfico, en dos columnas:
\begin{lstlisting}
\textbf{Problema.} Hallar el vértice de $f(x)=x^2-4x+1$.
\begin{columns}
  \begin{column}{0.55\textwidth}
    \begin{align*}
      f(x) &= x^2-4x+1 \\
           &= (x-2)^2 - 4 + 1 \\
           &= (x-2)^2 - 3.
    \end{align*}
    \begin{block}{Conclusión}
      Vértice $(2,-3)$; como $a>0$, mínimo $-3$ en $x=2$.
    \end{block}
  \end{column}
  \begin{column}{0.42\textwidth}
    \begin{tikzpicture}
    \begin{axis}[axis lines=middle, xlabel=$x$,
                 width=\linewidth, height=4cm,
                 xtick={2}, ytick={-3}]
      \addplot[blue,very thick,domain=-0.5:4.5]{x^2-4*x+1};
      \addplot[red,only marks,mark=*]
        coordinates {(2,-3)};
    \end{axis}
    \end{tikzpicture}
  \end{column}
\end{columns}
\end{lstlisting}
\end{frame}

\begin{frame}{Un problema resuelto: el resultado}
  \footnotesize
  \textbf{Problema.} Hallar el vértice de $f(x)=x^2-4x+1$.
  \begin{columns}
    \begin{column}{0.55\textwidth}
      \begin{align*}
        f(x) &= x^2-4x+1 \\
             &= (x-2)^2 - 4 + 1 \\
             &= (x-2)^2 - 3.
      \end{align*}
      \begin{block}{Conclusión}
        Vértice $(2,-3)$; como $a>0$, mínimo $-3$ en $x=2$.
      \end{block}
    \end{column}
    \begin{column}{0.42\textwidth}
      \begin{tikzpicture}
      \begin{axis}[axis lines=middle, xlabel=$x$,
                   width=\linewidth, height=4cm,
                   xtick={2}, ytick={-3}]
        \addplot[blue,very thick,domain=-0.5:4.5]{x^2-4*x+1};
        \addplot[red,only marks,mark=*]
          coordinates {(2,-3)};
      \end{axis}
      \end{tikzpicture}
    \end{column}
  \end{columns}
\end{frame}

\begin{frame}[fragile]{Una lámina-resumen: el código}
  \footnotesize Tres bloques de distinto color y un gráfico, en dos columnas:
\begin{lstlisting}
\begin{columns}
  \begin{column}{0.5\textwidth}
    \begin{block}{Definición}
      $f(x)=ax^2+bx+c$, con $a\neq0$.
    \end{block}
    \begin{alertblock}{Vértice y extremo}
      $x=-\frac{b}{2a}$; máx si $a<0$, mín si $a>0$.
    \end{alertblock}
    \begin{exampleblock}{Raíces}
      $x=\frac{-b\pm\sqrt{b^2-4ac}}{2a}$
    \end{exampleblock}
  \end{column}
  \begin{column}{0.46\textwidth}
    \begin{tikzpicture}
    \begin{axis}[axis lines=middle, width=\linewidth, height=4cm]
      \addplot[blue,very thick,domain=-1.2:3.2]{x^2-2*x-3};
      \addplot[red,only marks,mark=*]
        coordinates {(1,-4)};
      \addplot[orange,only marks,mark=*]
        coordinates {(-1,0) (3,0)};
    \end{axis}
    \end{tikzpicture}
  \end{column}
\end{columns}
\end{lstlisting}
\end{frame}

\begin{frame}{Una lámina-resumen: el resultado}
  \footnotesize
  \begin{columns}
    \begin{column}{0.5\textwidth}
      \begin{block}{Definición}
        $f(x)=ax^2+bx+c$, con $a\neq0$.
      \end{block}
      \begin{alertblock}{Vértice y extremo}
        $x=-\frac{b}{2a}$; máx si $a<0$, mín si $a>0$.
      \end{alertblock}
      \begin{exampleblock}{Raíces}
        $x=\frac{-b\pm\sqrt{b^2-4ac}}{2a}$
      \end{exampleblock}
    \end{column}
    \begin{column}{0.46\textwidth}
      \begin{tikzpicture}
      \begin{axis}[axis lines=middle, width=\linewidth, height=4cm]
        \addplot[blue,very thick,domain=-1.2:3.2]{x^2-2*x-3};
        \addplot[red,only marks,mark=*]
          coordinates {(1,-4)};
        \addplot[orange,only marks,mark=*]
          coordinates {(-1,0) (3,0)};
      \end{axis}
      \end{tikzpicture}
    \end{column}
  \end{columns}
\end{frame}

% ####################################################################
\section{Cierre}
% ####################################################################
\begin{frame}{Síntesis: la receta completa}
  \begin{enumerate}
    \item Elegí la \textbf{clase}: \texttt{article} para documentos,
      \texttt{beamer} para presentaciones.
    \item Armá el \textbf{preámbulo} con los \texttt{\textbackslash usepackage}
      que necesites (\texttt{babel}, \texttt{amsmath}, \texttt{graphicx},
      \texttt{pgfplots}\dots).
    \item Escribí el \textbf{cuerpo}: texto, \texttt{frame}s, bloques,
      columnas, fórmulas, tablas, imágenes y gráficos.
    \item \textbf{Compilá} con \texttt{latexmk -pdf} y revisá el PDF.
  \end{enumerate}
  \vfill
  \begin{block}{Para seguir aprendiendo}
    \begin{itemize}
      \item \texttt{Overleaf} (editor online con plantillas y tutoriales).
      \item La guía ``\emph{The Not So Short Introduction to \LaTeX}''.
      \item \texttt{es.wikibooks.org/wiki/Manual\_de\_LaTeX}.
    \end{itemize}
  \end{block}
\end{frame}

\begin{frame}[plain]
  \centering
  \vfill
  {\Huge \structure{¡Gracias!}}\\[1em]
  {\large Cómo crear documentos y presentaciones con LaTeX}\\[0.5em]
  Rodrigo Martín Vidal Galán
  \vfill
\end{frame}

\end{document}
